% GENERATED FILE. Edit canonical lessons, then run npm run generate:books.
% canonical-source-hash: fnv1a64:4da5dd51e6a22c92
% canonical-lessons: LA-C293-paulatim, LA-C293-sensim, LA-C293-gradatim, LA-C293-paulisper, LA-C293-aliquamdiu

\chapter{Paulātim}
\label{ch:la-paulatim}
\hlchaptermodality{293}

\begin{chapteropening}
\textbf{By the end of this chapter:} \emph{I can tell how slowly things happened: little by little, gradually, step by step, for a little while and for some time.}

Complete the last lesson of chapter 293: I can tell how slowly things happened: little by little, gradually, step by step, for a little while and for some time.
\end{chapteropening}

\section[paulātim]{paulātim — little by little}

\label{lesson:LA-C293-paulatim}



\begin{quote}
Before the new one: say the Latin for in an instant, then the Latin for immediately.
\end{quote}

\subsection*{You'll want to know: paulātim}
\textbf{paulātim} — \textquotedblleft{}little by little\textquotedblright{}.

\begin{quote}
Paulātim aqua in fossā crēscēbat.
\end{quote}

\begin{quote}
\emph{Little by little the water in the ditch was rising.}
\end{quote}

Built on \emph{paulum}, 'a little'. Caesar writes of people who \emph{paulātim} get used to something.

\begin{cousinweb}
\textbf{paulātim} is built on \emph{paulum}, \textquotedblleft{}a little\textquotedblright{} (Lewis \& Short), with the adverb ending \emph{-ātim} that is also seen in \emph{gradātim}: one little bit after another.
\end{cousinweb}

\begin{grammarlens}[title={What you've built}]
One more word for telling, explaining and describing at B1.
\end{grammarlens}

\subsection*{Your turn}
\begin{itemize}
\raggedright
  \item \emph{Say it:} \emph{paulātim}
  \item \emph{Say it:} \emph{paulātim}, once more
  \item \emph{Say it:} a sentence of your own with \emph{paulātim}
  \item \emph{Recall:} say the Latin for in an instant, then the Latin for immediately, then say \emph{paulātim} again
\end{itemize}

\begin{tcolorbox}[breakable,colback=teal!4,colframe=teal!35!black,title={Before you move on}]
What does paulātim mean? (\textquotedblleft{}Little by little\textquotedblright{}.) Say it once more.
\end{tcolorbox}
\section[sēnsim]{sēnsim — gradually, gently}

\label{lesson:LA-C293-sensim}



\begin{quote}
Before the new one: say the Latin for immediately, then the Latin for little by little.
\end{quote}

\subsection*{You'll want to know: sēnsim}
\textbf{sēnsim} — \textquotedblleft{}gradually, gently\textquotedblright{}.

\begin{quote}
Senex sēnsim ad forum ambulābat.
\end{quote}

\begin{quote}
\emph{The old man was walking slowly to the forum.}
\end{quote}

So gentle that each step is felt; the opposite of \emph{repente}.

\begin{cousinweb}
\textbf{sēnsim} is built on \emph{sentiō}, \textquotedblleft{}to perceive\textquotedblright{} (Lewis \& Short): first \textquotedblleft{}perceptibly\textquotedblright{}, then \textquotedblleft{}gently, little by little\textquotedblright{}, as opposed to what comes suddenly.
\end{cousinweb}

\begin{grammarlens}[title={What you've built}]
One more word for telling, explaining and describing at B1.
\end{grammarlens}

\subsection*{Your turn}
\begin{itemize}
\raggedright
  \item \emph{Say it:} \emph{sēnsim}
  \item \emph{Say it:} \emph{sēnsim}, once more
  \item \emph{Say it:} a sentence of your own with \emph{sēnsim}
  \item \emph{Recall:} say the Latin for immediately, then the Latin for little by little, then say \emph{sēnsim} again
\end{itemize}

\begin{tcolorbox}[breakable,colback=teal!4,colframe=teal!35!black,title={Before you move on}]
What does sēnsim mean? (\textquotedblleft{}Gradually, gently\textquotedblright{}.) Say it once more.
\end{tcolorbox}
\section[gradātim]{gradātim — step by step}

\label{lesson:LA-C293-gradatim}



\begin{quote}
Before the new one: say the Latin for little by little, then the Latin for gradually.
\end{quote}

\subsection*{You'll want to know: gradātim}
\textbf{gradātim} — \textquotedblleft{}step by step\textquotedblright{}.

\begin{quote}
Gradātim ad summum montem ascendimus.
\end{quote}

\begin{quote}
\emph{Step by step we climbed to the top of the mountain.}
\end{quote}

From \emph{gradus}, 'a step' (taught). Cicero uses it of an argument that moves one point at a time.

\begin{cousinweb}
\textbf{gradātim} is built on \emph{gradus}, \textquotedblleft{}step\textquotedblright{} (EtymDB), with the adverb ending \emph{-ātim}. \emph{Gradus} is inherited and goes back to a root meaning \textquotedblleft{}to walk, to step\textquotedblright{}.
\end{cousinweb}

\begin{grammarlens}[title={What you've built}]
One more word for telling, explaining and describing at B1.
\end{grammarlens}

\subsection*{Your turn}
\begin{itemize}
\raggedright
  \item \emph{Say it:} \emph{gradātim}
  \item \emph{Say it:} \emph{gradātim}, once more
  \item \emph{Say it:} a sentence of your own with \emph{gradātim}
  \item \emph{Recall:} say the Latin for little by little, then the Latin for gradually, then say \emph{gradātim} again
\end{itemize}

\begin{tcolorbox}[breakable,colback=teal!4,colframe=teal!35!black,title={Before you move on}]
What does gradātim mean? (\textquotedblleft{}Step by step\textquotedblright{}.) Say it once more.
\end{tcolorbox}
\section[paulisper]{paulisper — for a little while}

\label{lesson:LA-C293-paulisper}



\begin{quote}
Before the new one: say the Latin for gradually, then the Latin for step by step.
\end{quote}

\subsection*{You'll want to know: paulisper}
\textbf{paulisper} — \textquotedblleft{}for a little while\textquotedblright{}.

\begin{quote}
Paulisper in tabernā mānsimus; deinde discessimus.
\end{quote}

\begin{quote}
\emph{We stayed in the shop for a little while; then we left.}
\end{quote}

It says how long, briefly. In Plautus a slave begs his master to wait \emph{paulisper}.

\begin{cousinweb}
\textbf{paulisper} is built on \emph{paulum}, \textquotedblleft{}a little\textquotedblright{}, the neuter of \emph{paulus}, \textquotedblleft{}small\textquotedblright{}, and \emph{per}, \textquotedblleft{}through\textquotedblright{} (EtymDB, Lewis \& Short): through a small stretch of time.
\end{cousinweb}

\begin{grammarlens}[title={What you've built}]
One more word for telling, explaining and describing at B1.
\end{grammarlens}

\subsection*{Your turn}
\begin{itemize}
\raggedright
  \item \emph{Say it:} \emph{paulisper}
  \item \emph{Say it:} \emph{paulisper}, once more
  \item \emph{Say it:} a sentence of your own with \emph{paulisper}
  \item \emph{Recall:} say the Latin for gradually, then the Latin for step by step, then say \emph{paulisper} again
\end{itemize}

\begin{tcolorbox}[breakable,colback=teal!4,colframe=teal!35!black,title={Before you move on}]
What does paulisper mean? (\textquotedblleft{}For a little while\textquotedblright{}.) Say it once more.
\end{tcolorbox}
\section[aliquamdiū]{aliquamdiū — for some time}

\label{lesson:LA-C293-aliquamdiu}



\begin{quote}
Before the new one: say the Latin for step by step, then the Latin for for a little while.
\end{quote}

\subsection*{You'll want to know: aliquamdiū}
\textbf{aliquamdiū} — \textquotedblleft{}for some time\textquotedblright{}.

\begin{quote}
Aliquamdiū pugnātum est; tum hostēs fūgērunt.
\end{quote}

\begin{quote}
\emph{There was fighting for some time; then the enemy fled.}
\end{quote}

\emph{aliquam} + \emph{diū} (taught): for a fair while, longer than \emph{paulisper}. Caesar uses it of battles.

\begin{cousinweb}
\textbf{aliquamdiū} joins a form of \emph{aliquis}, \textquotedblleft{}some\textquotedblright{}, with \emph{diū}, \textquotedblleft{}for a long time\textquotedblright{} (Lewis \& Short): for some length of time.
\end{cousinweb}

\begin{grammarlens}[title={What you've built}]
One more word for telling, explaining and describing at B1.
\end{grammarlens}

\subsection*{Your turn}
\begin{itemize}
\raggedright
  \item \emph{Say it:} \emph{aliquamdiū}
  \item \emph{Say it:} \emph{aliquamdiū}, once more
  \item \emph{Say it:} a sentence of your own with \emph{aliquamdiū}
  \item \emph{Recall:} say the Latin for step by step, then the Latin for for a little while, then say \emph{aliquamdiū} again
\end{itemize}

\begin{tcolorbox}[breakable,colback=teal!4,colframe=teal!35!black,title={Before you move on}]
What does aliquamdiū mean? (\textquotedblleft{}For some time\textquotedblright{}.) Say it once more.
\end{tcolorbox}
